% !PRINTED 585
\weberpage{585}
%<<CONT
therefore takes only the product of all the mutually distinct ones among the conjugate ideals, this need by no means be an ideal in $R$.

What has been set forth here for the norms holds also for the other symmetric functions, in particular for the traces. We call the sum
\weq{6}{\mathfrak{S}_R(\omega) = \omega_1 + \omega_2 + \cdots + \omega_n}
the \gesperrt{partial trace (or relative trace) of $\omega$ with respect to} $R$. It is a number or a functional in $R$, and for the total trace we obtain
\weq{7}{S_\Omega(\omega) = S_R\,\mathfrak{S}_R(\omega),}
in which the meaning of the symbols is immediately intelligible.

A particularly important relation results, however, from the consideration of the fundamental ideals.

Let $\tau$ be a basis form of $\Omega$ (thus a linear form in $mn$ variables, \S.\,146), and let
\[
F(t) = N_\Omega(t - \tau)
\]
be the irreducible function of degree $mn$ (in the absolute domain of rationality) whose root is $\tau$. Then in \S.\,158 we designated the ideal defined by the functional $F'(\tau)$ as the \gesperrt{fundamental ideal} of the field $\Omega$.

If now $\sigma$ is a basis form of $R$ and
\[
\psi(t) = N_R(t - \sigma),
\]
then the fundamental ideal of the field $R$ is defined by $\psi'(\sigma)$.

From the function $F(t)$, however, there splits off in the field $R$ an irreducible factor of degree $n$,
\[
f(t) = \mathfrak{N}_R(t - \tau),
\]
which vanishes for $t = \tau$. By $f'(\tau)$ we define the \gesperrt{partial fundamental ideal of $\Omega$ with respect to} $R$, and now prove the theorem:
\weq{8}{F'(\tau) = \varepsilon\, f'(\tau)\, \psi'(\sigma),}
in which $\varepsilon$ is a unit.

To prove it, we set, following the notation employed in \S.\,158,
\weq{9}{\frac{F(t)}{t - \tau} = t^{mn-1}F_0(\tau) + t^{mn-2}F_1(\tau) + \cdots + F_{mn-1}(\tau),}
\weq{10}{\frac{f(t)}{t - \tau} = t^{n-1}f_0(\tau) + t^{n-2}f_1(\tau) + \cdots + f_{n-1}(\tau),}
\weq{11}{\frac{\psi(t)}{t - \sigma} = t^{m-1}\psi_0(\sigma) + t^{m-2}\psi_1(\sigma) + \cdots + \psi_{m-1}(\sigma).}
